\chapter{Discussão dos resultados}
\label{cap:discussao}

A discussão dos resultados tem como objetivo interpretar os dados obtidos durante a execução dos algoritmos e relacioná-los aos conceitos de arquitetura de computadores e organização da memória. A comparação entre os percursos linha por linha e coluna por coluna permite observar que a ordem de acesso aos dados influencia diretamente o desempenho de aplicações que trabalham com matrizes.

\begin{table}[H]
    \centering
    \caption{Comparação entre os dois padrões de acesso à memória}
    \label{tab:comparacao}
    \begin{tabular}{lcc}
        \toprule
        \textbf{Métrica} & \textbf{Linha por linha} & \textbf{Coluna por coluna} \\
        \midrule
        Tempo de execução (s) & 0,006974 & 0,010819 \\
        Soma calculada & 4194304 & 4194304 \\
        L1 cache misses & 61.204 & 5.444.797 \\
        LLC cache misses & 126.655 & 320.348 \\
        \bottomrule
    \end{tabular}
    \fonte{Elaborado pelos autores com dados do \texttt{perf stat} (2026).}
\end{table}

\section{Análise da comparação}
Os resultados demonstram uma diferença significativa de desempenho entre os dois métodos. O percurso linha por linha executou a tarefa em aproximadamente 0,006 segundos, enquanto o percurso coluna por coluna levou cerca de 0,011 segundos. Isso representa um tempo cerca de 55\% maior para o algoritmo que percorre as colunas.

\section{Do código em C ao hardware}
A execução de um programa em C envolve uma interação direta entre o software e os componentes físicos. Quando um programa acessa um elemento de uma matriz, o processador precisa obter o dado na hierarquia de memória. Na linguagem C, a organização \textit{row-major} faz com que o percurso linha por linha aproveite melhor os dados já carregados na cache.

No percurso linha por linha, o processador acessa posições consecutivas da memória, favorecendo o aproveitamento dos blocos de dados. No percurso coluna por coluna, os acessos ocorrem em posições mais distantes, provocando mais falhas de cache. Dessa forma, mesmo com a mesma complexidade assintótica (\(O(n \times m)\)), os desempenhos práticos são diferentes.

\section{Limitações do método}
Os resultados apresentados possuem limitações, como a dependência das configurações do compilador e as características do processador. Variações nos tempos de execução podem ocorrer devido ao sistema operacional e à execução de outros processos.

\section{Limitações físicas}
As limitações físicas estão relacionadas às características dos componentes de hardware. A memória cache possui capacidade limitada, enquanto a memória RAM apresenta maior latência. A largura de banda da memória e o tamanho dos blocos transferidos também influenciam o desempenho.

\section{Impacto no desenvolvimento de software de alto desempenho}
O estudo demonstra que a organização dos dados e a escolha dos padrões de acesso à memória são fundamentais. A utilização adequada da memória cache pode reduzir o tempo de execução e melhorar a eficiência dos recursos computacionais, especialmente em aplicações científicas, processamento de imagens e inteligência artificial.